\section{\label{sec:mixers-as-walks}Constraint-preserving mixers as
quantum walks}

Constraint-preserving mixers and quantum walks on feasible graphs are closely related, but the relationship is asserted case by case rather than established in general. Marsh and Wang~\cite{Marsh2020} introduced QWOA and noted that the transverse-field mixer is the adjacency matrix of a hypercube ~\cite{Marsh2019}. Slate et al.~\cite{Slate2021} apply QWOA to portfolio optimization with a walk on the complete graph over an indexed set of canonical solutions. Bärtschi and Eidenbenz~\cite{Bartschi2020} define the Grover mixer as a selective phase shift generated by the projector $\ket{s_\F}\bra{s_\F}$, without graph-theoretic language. Matwiejew and Wang~\cite{Matwiejew2024} analyse mixing unitaries defined by vertex-transitive graphs through automorphism orbits. Fuchs et al.~\cite{Fuchs2022,Fuchs2024} work in the complementary direction, building mixers that preserve a prescribed subspace.

What is missing is a statement covering an arbitrary feasibility-preserving mixer.

{\color{red}\textbf{[REVIEW -- scope/novelty]} This overstates the result. The lemma covers a restricted matrix class. Frame novelty around the unified dictionary and its consequences, or formally extend the framework to signed/complex hopping and on-site potentials.} Each identification above is settled by inspecting one particular Pauli construction. That the Grover mixer and the complete-graph walk of Ref.~\cite{Slate2021} are the same operator, since
\begin{equation}
    \mathcal{A}_{K_{\F}} = |\F|\,\ket{s_{\F}}\bra{s_{\F}} - I_{\F} ,
    \label{eq:grover-complete}
\end{equation}
is stated in Ref.~\cite{Bridi2024}, without the explicit rescaling, but in neither of the two references it relates. Lemma~\ref{lem:dictionary} settles such questions once.

\subsection{\label{subsec:dictionary}When is a mixer a quantum walk?}

\begin{lemma}[Mixer--walk correspondence]
\label{lem:dictionary}
Let $H_M$ be Hermitian with $H_M\mathcal{H}_{\F}\subseteq\mathcal{H}_{\F}$, and let $A_{xy}=\bra{x}H_M\ket{y}$ for $x,y\in\F$. Assume
\begin{enumerate}
    \item[(i)] the diagonal is constant, $A_{xx}=c$;
    \item[(ii)] the nonzero off-diagonal elements are real and of a common sign
$\varepsilon=\pm1$.
\end{enumerate}
Define the weighted feasible graph $G_{\F}=(\F,E_{\F})$ by
\begin{equation}
    E_{\F}=\bigl\{\{x,y\}:x\neq y,\ A_{xy}\neq 0\bigr\},
    \label{eq:dictionary-graph}
\end{equation}
with weights $w_{xy}=\varepsilon A_{xy}>0$. Then
\begin{equation}
    H_M\big|_{\mathcal{H}_{\F}}
    =
    \varepsilon\,\mathcal{A}_{\F}+c\,I_{\F},
    \label{eq:dictionary-operator}
\end{equation}
and consequently
\begin{equation}
    e^{-i\beta H_M}\big|_{\mathcal{H}_{\F}}
    =
    e^{-i\beta c}\,
    e^{-i(\varepsilon\beta)\mathcal{A}_{\F}} .
    \label{eq:dictionary-unitary}
\end{equation}
If in addition $\ket{\psi_0}=\ket{s_{\F}}$, the QAOA-type ansatz \eqref{eq:qaoa-type-state} is identical to the QWOA state \eqref{eq:qwoa-state} on $(G_{\F},w)$, up to a global phase and the reparametrization $\beta_\ell\mapsto\varepsilon\beta_\ell$.
\end{lemma}

\begin{proof}
Since $H_M$ preserves $\mathcal{H}_{\F}$, its restriction is determined by the block $(A_{xy})_{x,y\in\F}$. Splitting that block into diagonal and off-diagonal parts and using (i) and (ii),
\begin{align}
    H_M\big|_{\mathcal{H}_{\F}}
    &=
    c\,I_{\F}
    +
    \!\!\sum_{\{x,y\}\in E_{\F}}\!\!
    A_{xy}\bigl(\ket{x}\bra{y}+\ket{y}\bra{x}\bigr)
    \nonumber\\
    &=
    c\,I_{\F}+\varepsilon\,\mathcal{A}_{\F},
\end{align}
Hermiticity giving $A_{yx}=A_{xy}^{*}=A_{xy}$ for real entries. Equation~\eqref{eq:dictionary-unitary} follows because $I_{\F}$ commutes with $\mathcal{A}_{\F}$, so the exponential factorizes and the identity piece contributes a global phase. Global phases cancel in all measurement probabilities, and $\beta_\ell\mapsto\varepsilon\beta_\ell$ is a bijection of the parameter domain, so the two ans\"atze have identical optimization landscapes.

{\color{red}\textbf{[REVIEW -- parameter-domain caveat]} Only if the allowed parameter domains are mapped bijectively under the stated reparametrization. State the numerical domains explicitly.}
\end{proof}

The content of the lemma is in its hypotheses: given (i) and (ii), Eq.~\eqref{eq:dictionary-operator} is the definition of a weighted graph read backwards. What it supplies is a check that is uniform across mixers, and that is complete in the sense of the following remark. It assumes feasibility preservation rather than establishing it; constructing mixers that preserve a prescribed subspace is a separate problem~\cite{Fuchs2022,Fuchs2024}.

\begin{remark}  
\label{rem:failure}
If the diagonal is not constant, write $A_{xx}=V(x)$. Then
\begin{equation}
    H_M\big|_{\mathcal{H}_{\F}}
    =
    \mathcal{A}_{\F}+\sum_{x\in\F}V(x)\ket{x}\bra{x},
    \label{eq:potential}
\end{equation}
so the mixer generates a walk with an additional on-site potential term rather than a pure walk, and $\ket{s_{\F}}$ is in general not an eigenvector even when $G_{\F}$ is regular. Penalty formulations fall into this class~\cite{Lucas2014}, since the penalty contributes a configuration-dependent diagonal.

{\color{red}\textbf{[REVIEW -- conceptual error]} Standard penalty-QAOA adds the penalty to $H_C$, not to $H_M$. It therefore does not turn the mixer into a walk with an on-site potential. Rewrite this distinction explicitly.} Note that $V$ cannot be absorbed into the phase separator, as the two enter the ansatz at different points in the ordered product. If instead the off-diagonal elements are complex, $A_{xy}=|A_{xy}|e^{i\phi_{xy}}$, the phases can be removed by a diagonal unitary only when they sum to zero modulo $2\pi$ around every cycle of $G_{\F}$; otherwise the evolution differs from that of any real-weighted walk on $G_{\F}$. All mixers below have constant diagonal and real non-negative off-diagonal elements, and fall under Lemma~\ref{lem:dictionary} directly.
\end{remark}

\subsection{\label{subsec:transverse}Transverse-field mixer}

That the transverse-field mixer generates a walk on the hypercube is standard~\cite{Marsh2019,Marsh2020,Slate2021,Matwiejew2024} and we record it for reference. Since $X_i\ket{x}=\ket{x\oplus e_i}$ and $y=x\oplus e_i$ holds for at most one index $i$,
\begin{equation}
    \bra{y}\sum_{i=1}^{n}X_i\ket{x}
    =
    \begin{cases}
        1, & d_H(x,y)=1,\\
        0, & \text{otherwise},
    \end{cases}
    \label{eq:hypercube}
\end{equation}
so that $\sum_i X_i=\mathcal{A}_{Q_n}$. As $\ket{+}^{\otimes n}=\ket{s_{\{0,1\}^n}}$, standard QAOA~\cite{Farhi2014} is exactly QWOA on $Q_n$, with no rescaling and no residual phase. The sign convention $H_M=-\sum_i X_i$ is the case $\varepsilon=-1$ of Lemma~\ref{lem:dictionary}. The hypercube walk does not preserve the cardinality constraint, which is why penalty terms are required in the unconstrained formulation~\cite{Lucas2014,Hodson2019}; by Remark~\ref{rem:failure} the penalized mixer then carries an on-site potential.

\subsection{\label{subsec:xy}$XY$ mixers}

Constraint-preserving mixers for fixed Hamming weight are built from the two-qubit $XY$ interaction~\cite{Wang2020}
\begin{equation}
    H^{(ij)}_{XY}
    =
    \tfrac{1}{2}\left(X_iX_j+Y_iY_j\right),
    \label{eq:xy-pair}
\end{equation}
which, given a connectivity graph $G=([n],E)$ on the qubits (i.e., a graph whose vertices are the $n$ qubit labels and an edge $\{i,j\}\in E$ declares that qubits $i$ and $j$ are allowed to exchange their values), defines
\begin{equation}
    H_{XY}(G)
    =
    \sum_{\{i,j\}\in E}
    H^{(ij)}_{XY} .
    \label{eq:xy-mixer}
\end{equation}
Note that $G$ lives on the $n$ qubits, whereas the feasible graph $G_\F$ lives on the $\binom{n}{k}$ feasible configurations. Some authors~\cite{Wang2020,Matwiejew2024} omit the factor $\tfrac12$ in Eq.~\eqref{eq:xy-pair}, in which case all statements below hold with $\beta\mapsto2\beta$.

\begin{proposition}[$XY$ mixers are quantum walks]
\label{prop:xy-hopping}
Let $G=([n],E)$ and $1\le k\le n-1$, and let $G_{\F}[G]$ be the graph with vertex set $\F_{n,k}$ in which
\begin{equation}
    x\sim y
    \iff
    \operatorname{supp}(x)\,\triangle\,\operatorname{supp}(y)\in E ,
    \label{eq:induced-adjacency}
\end{equation}
$\triangle$ denoting symmetric difference; equivalently, $y$ is obtained from $x$ by moving one occupied site along an edge of $G$. Then
\begin{equation}
    H_{XY}(G)\big|_{\mathcal{H}_{\F_{n,k}}}
    =
    \mathcal{A}_{G_{\F}[G]} .
    \label{eq:xy-hopping}
\end{equation}
\end{proposition}

\begin{proof}
In the two-qubit computational basis the weight-changing terms of $X_iX_j+Y_iY_j$ cancel and the weight-preserving terms add, giving
\begin{equation}
    H^{(ij)}_{XY}
    =
    \ket{0_i1_j}\bra{1_i0_j}+\ket{1_i0_j}\bra{0_i1_j}
    \label{eq:xy-hopping-form}
\end{equation}
tensored with the identity elsewhere; in particular $H^{(ij)}_{XY}$ annihilates $\ket{0_i0_j}$ and $\ket{1_i1_j}$. Therefore $H_{XY}(G)$ preserves every fixed-weight subspace and satisfies Eq.~\eqref{eq:feasible-condition}. Let $x\in\F_{n,k}$ with support $S=\operatorname{supp}(x)$. By Eq.~\eqref{eq:xy-hopping-form}, the term $H^{(ij)}_{XY}$ acts non-trivially on $\ket{x}$ precisely when exactly one of $i,j$ lies in $S$; in that case $H^{(ij)}_{XY}\ket{x}=\ket{y}$, where $y$ is obtained from $x$ by moving the single occupied site from one endpoint of the edge to the other. Writing $S'=\operatorname{supp}(y)$, this says $|S|=|S'|=k$ and $S\,\triangle\,S'=\{i,j\}$.

Conversely, if $y\in\F_{n,k}$ satisfies $S\,\triangle\,S'=\{i,j\}$ for some $\{i,j\}\in E$, then exactly one of $i,j$ lies in $S$ and $y$ arises from the corresponding edge term. Thus the strings reachable from $x$ are exactly those whose supports differ from $S$ by a symmetric difference that is an edge of $G$, each reached by a single edge term, so
\begin{equation}
    H_{XY}(G)\ket{x}=\sum_{y\sim x}\ket{y},
\end{equation}
with $\sim$ as in Eq.~\eqref{eq:induced-adjacency}. No edge term contributes a $\ket{x}\bra{x}$ component, so the diagonal vanishes. Applying Lemma~\ref{lem:dictionary} with $c=0$, $\varepsilon=+1$ and unit weights completes the proof.
\end{proof}

The graphs $G_{\mathcal{F}}[G]$ are isomorphic to a class of graphs in the literature known as token graphs, whose general theory is developed in Refs.~\cite{FabilaMonroy2012Token,Carballosa2017Regularity}. Proposition~\ref{prop:xy-regularity} corresponds to a result established in Ref.~\cite{Carballosa2017Regularity}.

{\color{red}\textbf{[REVIEW -- novelty boundary]} Make explicit that the graph-theoretic regularity theorem is prior work; the new contribution is its interpretation/specialization for $XY$ mixers and algorithm design.}

{\color{red}\textbf{[LITERATURE UPDATE -- XY topology]} The recent DLA analysis of Kordonowy and Leipold~\cite{Kordonowy2026} explicitly studies different $XY$-mixer topologies for cardinality-constrained optimization. This paper should be discussed here because it changes the novelty boundary: topology-dependent $XY$ expressivity is now a direct literature topic. The present result remains distinct in identifying the induced feasible-state/token graph and embedding complete-$XY$ into the Johnson association scheme, but claims about topology, expressivity, or trainability should be compared against Ref.~\cite{Kordonowy2026}.}

\begin{corollary}[Complete $XY$ mixer]
\label{cor:dicke-xy-johnson}
For $G=K_n$,
\begin{equation}
\begin{gathered}
    H_{XY}(K_n)\big|_{\mathcal{H}_{\F}}=\mathcal{A}_{J(n,k)},\\
    \ket{D^n_k}=\ket{s_{\F}} .
\end{gathered}
\end{equation}
The Dicke-state plus complete-$XY$ ansatz is therefore exactly QWOA on $J(n,k)$, with no global phase and no rescaling of $\beta$.
\end{corollary}

{\color{red}\textbf{[NOVELTY AUDIT 2026 -- Johnson identity is prior art]} This operator identity is correct, but it is no longer defensible as a standalone novelty claim. Structurally, the $k$-token graph of $K_n$ is the Johnson graph in the token-graph literature~\cite{FabilaMonroy2012Token}; algorithmically, Min, Seo, and Heo~\cite{MinSeoHeo2026} explicitly state that an $XY$ mixer confines dynamics to Johnson-graph subspaces. Recast this corollary as part of the paper's unifying dictionary and use it to motivate the Johnson-scheme extension, which is the more distinctive contribution.}

{\color{red}\textbf{[NOVELTY AUDIT 2026 -- prior XY topology studies]} Ring and complete $XY$ mixers and Dicke-state initialization were already compared in the quantum alternating-operator literature~\cite{Cook2019,Wang2020,He2023}. Kordonowy and Leipold~\cite{Kordonowy2026} further analyze the dynamical Lie algebras of $XY$-mixer topologies. Therefore the manuscript should avoid implying that topology dependence of $XY$ mixers itself is new; what is new must be tied to the induced feasible-state graph/Johnson-scheme family and the controlled sweep performed here.

\begin{proof}
Apply Proposition~\ref{prop:xy-hopping} with $E=\{\{i,j\}:i\neq j\}$: every pair is an edge, so $x\sim y$ iff $|\operatorname{supp}(x)\triangle\operatorname{supp}(y)|=2$, iff the $k$-subsets differ in exactly one element, iff $d_H(x,y)=2$, which is the definition of $J(n,k)$. Comparing Eq.~\eqref{eq:dicke-state} with Eq.~\eqref{eq:uniform-feasible-state} gives $\ket{D^n_k}=\ket{s_{\F}}$.
\end{proof}

\begin{remark} The Johnson graph is regular and vertex-transitive, so
$\ket{D^n_k}$ is its top eigenvector. For $k=1$ one recovers $G_{\F}[G]\cong G$, the single-particle walk on the connectivity graph.
\end{remark}

\begin{corollary}[Ring $XY$ mixer]
\label{cor:ring-xy}
Let $G=C_n$ be the cycle. Then $H_{XY}(C_n)|_{\mathcal{H}_{\F}}=\mathcal{A}_{G_{\F}[C_n]}$, where $G_{\F}[C_n]$ is a connected spanning subgraph of $J(n,k)$ with
\begin{equation}
    \deg(x)=2\,b(x),
    \label{eq:ring-degree}
\end{equation}
$b(x)$ being the number of maximal cyclic blocks of ones in $x$. In particular $G_{\F}[C_n]$ is not regular for $2\le k\le n-2$, and $\ket{D^n_k}$ is not an eigenvector of $H_{XY}(C_n)$.
\end{corollary}

{\color{red}\textbf{[LITERATURE UPDATE -- initial-state/mixer alignment]} He et al.~\cite{He2023} directly compare ring and complete $XY$ mixers for constrained portfolio optimization and show that alignment between the initial state and the mixer ground state affects performance; they also study Trotterized mixers and execute a 32-qubit trapped-ion experiment. This prior work should be cited explicitly when interpreting the Dicke/complete-$XY$ versus Dicke/ring-$XY$ distinction. Note that their ground-state convention must be reconciled with the sign convention used in this manuscript.}

\begin{proof}
The first statement is Proposition~\ref{prop:xy-hopping} with $E=E(C_n)$, and $G_{\F}[C_n]\subseteq J(n,k)$ follows from $E(C_n)\subset E(K_n)$. Given a vertex $x=(x_1,\dots,x_n)$ with $b(x)=m$ maximal cyclic blocks of ones, it can be written as $(0^{a_1}1^{b_1}\cdots 0^{a_m}1^{b_m})$, where $a_l,b_l\ge1$, $\sum_l b_l=k$ and $\sum_l(a_l+b_l)=n$. By Proposition~\ref{prop:xy-hopping}, the neighbours of $x$ are the vertices that differ in exactly two adjacent positions. Thus, for a given cyclic block of ones $1^{b_l}$ in this representation of $x$, $x$ possesses two neighbours: one swapping the leftmost $1$ with its adjacent $0$, and the other swapping the rightmost $1$ with its adjacent $0$. Consequently, each maximal cyclic block of ones yields exactly two neighbours, which implies that the degree of $x$ is $2b(x)$. For $2\le k\le n-2$ non-regularity follows because both values $b(x)=1$ and $b(x)=2$ occur: the first by taking the $k$ ones contiguous, which is possible since $k\le n-1$, and the second by splitting them into two blocks, which is possible since $k\ge2$ and $n-k\ge2$. Degrees $2$ and $4$ therefore both occur; for $n=4$, $k=2$ these are the strings $1100$ and $1010$. Connectivity holds since any occupied site can be moved to any position by successive nearest-neighbour hops.
\end{proof}

Corollary~\ref{cor:ring-xy} is not an isolated accident: regularity of $G_{\F}[G]$ is exceptional among connectivity graphs, and the exceptions can be listed.

\begin{proposition}[Regularity of $XY$ mixers]
\label{prop:xy-regularity}
Let $G=([n],E)$ and $1\le k\le n-1$. For $x\in\F_{n,k}$,
\begin{equation}
    \deg_{G_{\F}[G]}(x)
    =
    e_G\bigl(\operatorname{supp}(x),[n]\setminus\operatorname{supp}(x)\bigr),
    \label{eq:boundary-degree}
\end{equation}
the number of edges of $G$ with exactly one endpoint in $\operatorname{supp}(x)$. Hence $G_{\F}[G]$ is regular if and only if every $k$-subset of $[n]$ has the same edge boundary in $G$, and, for $n\ge4$:
\begin{enumerate}
\item[(i)] for $k=1$ and $k=n-1$, if and only if $G$ is regular;
\item[(ii)] for $2\le k\le n-2$ with $n\neq2k$, if and only if $G$ is
complete or empty;
\item[(iii)] for $n=2k$, if and only if $G$ is complete, empty, the star
$K_{1,n-1}$, or the complement of a star, the last two giving $G_{\F}[G]$ of degree $k$ and $k(k-1)$.
\end{enumerate}
\end{proposition}

For a detailed proof, see Refs.~\cite{FabilaMonroy2012Token,Carballosa2017Regularity}.

\begin{remark}
The star and its complement single out one asset, so they are not $S_n$-invariant and lie outside the scope of Proposition~\ref{prop:min-degree}; at $n=2k$ the sparsest $S_n$-invariant graph is instead the Kneser graph, of degree $v_{r_{\max}}=1$.
\end{remark}

The complete-$XY$ mixer walks on a vertex-transitive, distance-transitive graph and is therefore within the scope of the automorphism-orbit analysis of Ref.~\cite{Matwiejew2024}; the ring-$XY$ mixer is not. The parity graph tabulated in that reference is a different object, arising from the mixer split into two disjoint matchings over alternating qubit pairs rather than from the full cycle Hamiltonian. That $\ket{D^n_k}$ fails to be an eigenvector of $H_{XY}(C_n)$ is already known: Ref.~\cite{He2023} reports that QAOA performs best when the initial state is the ground state of the mixer, finds the exact ring-mixer ground state impractical to prepare, and uses the Dicke state as a proxy for it. Corollary~\ref{cor:ring-xy} provides a possible reason and quantifies it: the non-regularity of $G_\F[C_n]$ with the degree of Eq.~\eqref{eq:ring-degree}.

\subsection{\label{subsec:grover}Grover mixer}

The feasible Grover mixer is the reflection
\begin{equation}
    \mathcal{G}_{\F}
    =
    2\ket{s_{\F}}\bra{s_{\F}}-I_{\F},
    \label{eq:grover-reflection}
\end{equation}
with $I_{\F}$ the identity on $\mathcal{H}_{\F}$~\cite{Bartschi2020,Bridi2024}.

\begin{proposition}
\label{prop:grover-complete-qwoa}
For the fixed-cardinality feasible set $\F$, the Dicke-state plus Grover-mixer ansatz is equivalent to QWOA on the complete feasible graph $K_{\F}$, up to a global phase and a rescaling of the mixing parameter.
\end{proposition}

\begin{proof}
Let $M=|\F|=\binom{n}{k}$ and $J_{\F}=\sum_{x,y\in\F}\ket{x}\bra{y}$. Since $\ket{s_{\F}}\bra{s_{\F}}=J_{\F}/M$ and $J_{\F}=\mathcal{A}_{K_{\F}}+I_{\F}$,
\begin{equation}
    \mathcal{G}_{\F}
    =
    \frac{2}{M}\mathcal{A}_{K_{\F}}
    +\Bigl(\frac{2}{M}-1\Bigr)I_{\F},
\end{equation}
so that
\begin{equation}
    e^{-i\beta \mathcal{A}_{K_{\F}}}
    =
    e^{-i\beta\left(\frac{M}{2}-1\right)}
    e^{-i(M\beta/2)\mathcal{G}_{\F}},
    \label{eq:grover-rescaling}
\end{equation}
giving $\theta=M\beta/2$. This is an instance of Lemma~\ref{lem:dictionary}: $\mathcal{G}_\F$ has constant diagonal $2/M-1$ and constant positive off-diagonal $2/M$, so the associated graph is the complete graph on $\F$ with uniform weight $2/M$.
\end{proof}

The same algebra relates any union of distance classes to its complement within the family.

\begin{proposition}[Complementary unions]
\label{prop:complementary}
Let $\emptyset\neq T\subsetneq\{1,\dots,r_{\max}\}$ and $T^{c}=\{1,\dots,r_{\max}\}\setminus T$. Then
\begin{equation}
    \mathcal{A}_{A_T}+\mathcal{A}_{A_{T^{c}}}
    =
    \mathcal{A}_{K_{\F}}
    =
    M\ket{s_{\F}}\bra{s_{\F}}-I_{\F} ,
    \label{eq:complementary-sum}
\end{equation}
and, $\ket{s_{\F}}$ being an eigenvector of $\mathcal{A}_{A_T}$ by Proposition~\ref{prop:min-degree}(ii), the three terms commute and
\begin{equation}
    e^{-i\beta\mathcal{A}_{A_{T^{c}}}}
    =
    e^{i\beta}\,e^{+i\beta\mathcal{A}_{A_T}}
    \bigl[I_{\F}+(e^{-i\beta M}-1)\ket{s_{\F}}\bra{s_{\F}}\bigr] .
    \label{eq:complementary-unitary}
\end{equation}
The two QWOA ans\"atze therefore coincide up to a global phase and $\beta_\ell\mapsto-\beta_\ell$, apart from an extra phase $e^{-i\beta_\ell M}$ applied at each layer to the component of the state along $\ket{s_{\F}}$.
\end{proposition}

\begin{proof}
The classes $A_1,\dots,A_{r_{\max}}$ partition the pairs of distinct feasible configurations, so the edge sets of $A_T$ and $A_{T^{c}}$ partition that of $K_{\F}$ and the adjacency operators add to $\mathcal{A}_{K_{\F}}$, which is Eq.~\eqref{eq:grover-complete}. Since $\ket{s_{\F}}$ is an eigenvector of $\mathcal{A}_{A_T}$, the rank-one projector $\ket{s_{\F}}\bra{s_{\F}}$ commutes with it, and $I_{\F}$ commutes with both. Exponentiating the three commuting terms of $\mathcal{A}_{A_{T^{c}}}=M\ket{s_{\F}}\bra{s_{\F}}-I_{\F}-\mathcal{A}_{A_T}$ and using $e^{-i\beta M\ket{s_{\F}}\bra{s_{\F}}} =I_{\F}+(e^{-i\beta M}-1)\ket{s_{\F}}\bra{s_{\F}}$ gives Eq.~\eqref{eq:complementary-unitary}. Equation~\eqref{eq:grover-complete} is the case $T=\{1,\dots,r_{\max}\}$, where $\mathcal{A}_{A_{T^{c}}}=0$.
\end{proof}

The Dicke plus Grover ansatz thus realizes one specific QWOA mixer, and does not imply equivalence to QWOA on other feasible graphs such as $J(n,k)$. The same computation with $\F=\{0,1\}^n$ shows that the unconstrained Grover mixer is QWOA on $K_{2^n}$. The complete feasible graph is the mixer of the original QWOA~\cite{Marsh2019,Marsh2020,Slate2021} and has been realized experimentally for portfolio optimization~\cite{Qu2024}.

The identification transfers mathematical and computational properties between QAOA mixer constructions and quantum walks. Reference~\cite{Bridi2024} proves that the expectation value of QAOA with the Grover mixer is invariant under any permutation of the feasible states, so the algorithm depends on the cost function only through the multiset of its values, and derives an asymptotic bound limiting its advantage to that of unstructured search. By Proposition~\ref{prop:grover-complete-qwoa} the same holds for QWOA on $K_{\F}$. The Johnson graph, by contrast, records which configurations differ by a single exchange. This is what the complete feasible graph discards, and it predicts the ordering observed at the endpoints of the family in Sec.~\ref{sec:results}.

Table~\ref{tab:mixer-graph} collects the mixers treated in this section with the feasible graphs they generate.

\begin{table}[t]
\caption{\label{tab:mixer-graph}Constraint-preserving mixers and their
feasible graphs. The initial state is in each case the uniform superposition
over the feasible set, which is $\{0,1\}^n$ in the first two rows and
$\F_{n,k}$ in the rest. In the last row $G$ is an arbitrary connectivity
graph, whose mixer is regular only in the cases listed in
Proposition~\ref{prop:xy-regularity}.}
\begin{ruledtabular}
\setlength{\tabcolsep}{3pt}
\begin{tabular}{llll}
Mixer & Graph & Reg. & Ref. \\
\colrule
$\sum_i X_i$ & $Q_n$ & yes & \cite{Farhi2014,Marsh2019} \\
$2\ket{+^n}\bra{+^n}-I$ & $K_{2^n}$ & yes & \cite{Bartschi2020} \\
$H_{XY}(K_n)$ & $J(n,k)$ & yes & \cite{Wang2020,Matwiejew2024} \\
$H_{XY}(C_n)$ & $G_{\F}[C_n]$ & no & this work \\
$\mathcal{G}_{\F}$ & $K_{\F}$ & yes & \cite{Bartschi2020,Slate2021} \\
$H_{XY}(G)$ & $G_{\F}[G]$ & see Prop.~\ref{prop:xy-regularity} & this work \\
\end{tabular}
\end{ruledtabular}
\end{table}